The next lemma aids in the proof of Theorem (\ref{thm:stochasticDominance}).
\begin{lemma}\label{lemma:binLR}
Let $Z_n^{\eta}$ be a Binomial random variable with parameters $n\in \mathbb{N}$ and $\eta\in(0,1)$. Given a number of successes, $k\in\{0,\dots,n\}$, we know that the probability mass function of $Z_n^{\eta}$ is $f_k(\eta):=\Pr[Z_n^{\eta}=k]={n\choose k}\eta^k(1-\eta)^{n-k}$. Let $\mathcal{L}(\eta|k)$ be the likelihood function of the event $Z_n^{\eta}=k$. Then the maximum likelihood of $f_k(\eta)$ is $\eta=\frac{k}{n}$. I.e., 
\[
\mathcal{L}(\eta|k)=
\text{argmax}_{\eta}f_k(\eta)=
\text{argmax}_{\eta}{n\choose k}\eta^k(1-\eta)^{n-k}=\frac{k}{n}.
\]
\end{lemma}

\begin{proof}[Proof of Lemma \ref{lemma:binLR}]

We notice that ${n\choose k}$ does not depend on $\eta$, thus
\[
\text{argmax}_{\eta}f_k(\eta)=
\text{argmax}_{\eta}{n\choose k}\eta^k(1-\eta)^{n-k}=\text{argmax}_{\eta}\eta^k(1-\eta)^{n-k}\]
The log-likelihood is particularly convenient for maximum likelihood estimation. Logarithms are strictly increasing functions, as a result, maximizing the likelihood is equivalent to maximizing the log-likelihood, i.e.,
\[
\text{argmax}_{\eta}\eta^k(1-\eta)^{n-k}=\text{argmax}_{\eta}\ln(\eta^k(1-\eta)^{n-k})=\text{argmax}_{\eta}k\ln(\eta)+(n-k)\ln(1-\eta)
\]
Differentiating (with respect to $\eta$) and comparing to zero we get
\[
\frac{d\ln(f_k(\eta))}{d\eta}=\frac{k}{\eta}-\frac{n-k}{1-\eta}=0.
\]
And after refactoring,
\[
k(1-\eta)=(n-k)\eta
\]
The function $\ln(f_k(\eta))$ is a strictly concave as its second derivative is negative,
\[
\frac{d^2\ln(f_k(\eta))}{d\eta^2}=-\frac{k}{\eta^2}-\frac{n-k}{(1-\eta)^2}<0,
\]
And since the derivative of a strictly concave function is zero at $\frac{k}{n}$, then $\hat{\eta}=\frac{k}{n}$ is a global maximum.
Therefore, $\hat{\eta}=\frac{k}{n}$ obtains absolute maximum in $f_k(\eta)$.
\end{proof}


\begin{proof}[Proof of Theorem \ref{thm:stochasticDominance}]
Let $Z_\tau^{\eta_i}$ be a random variable 
that represents the number of flips out of a $\tau$-tests sequence 
with a noise level of $\eta_i$,
i.e., $Z_\tau^{\eta_i}$ is the number of times 
when $y_j\ne y$ for $1\leq j \leq \tau$. 
We use $Z_\tau^{\eta_i}$ to express $\Pr[\pi(\hat{y}_{i,1},\dots,\hat{y}_{i,\tau})=1|y_i=0,\eta=\eta_i]$ as the probability that at least $\theta$ flips,
\[
\Pr[\pi(\hat{y}_{i,1},\dots,\hat{y}_{i,\tau})=1|y_i=0,\eta=\eta_i]=
\Pr[Z_\tau^{\eta_i}\geq \theta]
\]
and the probability of $\Pr[\pi(\hat{y}_{i,1},\dots,\hat{y}_{i,\tau})=1|q=+1,\eta=\eta_i]$ as at most $\tau-\theta$ flips, thus 
\[
\Pr[\pi(\hat{y}_{i,1},\dots,\hat{y}_{i,\tau})=1|y_i=+1,\eta=\eta_i]=\Pr[Z_\tau^{\eta_i}\leq \tau- \theta].
\]
From Lemma (\ref{lemma:binLR}) and since  probability density function (pdf) are is monotone increasing, 
we derive that the pdf of $Z_n^{\eta_2}$ 
satisfies \emph{monotone likelihood ratio property} 
over the pdf of $Z_n^{\eta_1}$. 
This implies that the pdf of $Z_n^{\eta_2}$ 
also has \emph{first-order stochastic dominance} over $Z_n^{\eta_1}$ 
by Theorem 1.1 in \cite{Whitt1980Uniform}. 
From \emph{stochastic dominance}, 
we can derive the desired inequalities 
\begin{equation*}
FP_{\theta,\tau}^{\eta_1}=\Pr[\theta\leq Z_n^{\eta_1}]<\Pr[ \theta \leq Z_n^{\eta_2}]
=FP_{\theta,\tau}^{\eta_2}
\end{equation*}
and
\begin{equation*}
FN_{\theta,\tau}^{\eta_1}=\Pr[Z_n^{\eta_1}\leq \tau- \theta]<\Pr[Z_n^{\eta_2}\leq \tau- \theta]=FN_{\theta,\tau}^{\eta_2}.
\end{equation*}
\end{proof}